% !TEX root = ../../Main.tex
\section{Research Gap: Positioning This Work}

The eighteen studies agree that a deep reinforcement learning agent can learn a profitable trading policy. They differ, or say nothing, on the conditions in which this holds. The findings above describe what the field does. When they are read against one another, they also show what the field has not done yet. \Cref{Tables:Positioning} sets out the comparison.

% !TEX root = ../Main.tex
\begin{table}[htb!]
\centering
\footnotesize
\caption{How this work compares with the eighteen surveyed articles.}
\label{Tables:Positioning}
\begin{tabularx}{\textwidth}{@{}lXX@{}}
\toprule
\textbf{Dimension} & \textbf{Surveyed literature (18 articles)} & \textbf{This work} \\
\midrule
Market & Stocks 12, Forex 3, ETFs 2, Futures 1 & Forex, the seven major pairs \\
\addlinespace
Data resolution & Bars in 17 of 18, daily in 13; one study uses quotes & Quote-level tick data, with hourly and daily decisions \\
\addlinespace
Data span & Ten years or more in 12 of 18, median eleven years & Eleven years, 2015 to 2026 \\
\addlinespace
Action space & Discrete 15, continuous 3; one sizes the position & Continuous, sets direction and size together \\
\addlinespace
Reward function & Total return 14, Sharpe ratio 3, log return 1 & Log return; models selected by the Calmar ratio and a risk-adjusted composite score \\
\addlinespace
Trading costs & Charged in 5 of 18 & Quoted spread at the moment of the trade plus the commission of a swap-free raw-spread account \\
\addlinespace
Instruments per design & One agent per instrument, or one agent on a basket & One design trained separately on seven pairs, then combined \\
\addlinespace
Reported robustness & The candidate count and the sensitivity to starting conditions are not stated & Five starting balances per pair, with the candidate count stated \\
\bottomrule
\end{tabularx}
\end{table}

The first gap appears only when two of those dimensions are read together. The three studies that trade currencies \cite{carapuco_reinforcement_2018, huang_financial_2018, rundo_deep_2019} all limit the agent to a discrete choice. The three that give the agent a continuous action \cite{majidi_algorithmic_2022, jia_lstm-ddpg_2021, zhang_deep_2020} all trade something else. In the work reviewed here the two properties never appear together. So there is no direct comparison for a currency agent that sets its own position size. This work uses that combination, and for this reason it cannot be measured against an existing result.

The second gap comes from the cost figures. A strategy tested without costs is not tested against the market it claims to trade. In Forex this omission matters more than the reported returns suggest, because the spread is not fixed. It widens when liquidity falls and around releases of economic data. These are the moments when a model trained on an average spread would expect to act. The third gap is about resolution, not length of history. A long sample is already normal in this literature, so it does not set a study apart, but working inside the bar does. A bar reports four prices and hides the order in which the price reached them. That order decides which spread was paid, and whether a stop was reached before a target.

The fourth gap is about evidence, not design. Each study reports the performance of a model it selected. None of the studies reviewed here states how many candidates that model was selected from. None states how the reported figure changes when the evaluation is repeated from a different starting condition. Deep reinforcement learning is sensitive to initialisation. So a single figure cannot separate what comes from the method and what comes from the particular run. A model chosen as the best of many also tells as much about the sample it was chosen from as about the market it traded \cite{bailey_pseudo-mathematics_2014, lopez_de_prado_advances_2018}. The last gap is about how results are combined. The studies reviewed report either one agent trading one instrument, or one agent trading a basket. None of them trains one design separately on each of several instruments and then measures what those agents do when they run together. That is how a trader who holds several currency pairs would deploy them. The individual results do not show what the combination does, because that depends on how the equity curves move against each other.

This work is placed in these gaps. One design is applied without change to all seven major currency pairs. It uses a continuous action that sets the direction and the size of the position together. The evaluation runs on quote-level data over eleven years. It charges the spread quoted at the moment of each trade, plus the commission of a swap-free raw-spread retail account. This is an account that a retail trader can open, and not a simplification made for modelling. Each pair is reported from five different starting balances, not from a single backtest, and the number of candidates the selected model was drawn from is stated. The seven agents are then combined into an equal-weight portfolio. So this work reports both the seven single results and the one result that a trader running all seven would have seen. \Cref{Chapter:ExperimentalResults} gives both.

There is also one point where this work is the same as the literature. The reward is the log return. One of the surveyed studies makes the same choice \cite{majidi_algorithmic_2022}, so this is not a point of difference. The difference is that the reward the agent optimises and the criteria used to select the final model are, on purpose, not the same measure. During training, candidate runs are scored by the Calmar ratio. The model finally published for each pair is the one ranked highest by a composite score. This score is built from percentile ranks within the pair's own pool of candidates. It uses the annualised return, the Sharpe, Sortino, Calmar and Sterling ratios defined in \Cref{Section:PerformanceMeasurement}, the maximum and downside risk, and measures of how the policy behaves. So a policy that earns its return by taking more risk does not win on return alone. \Cref{Chapter:Implementation} sets out both.